% ============================================================
%  CFF Discovery: Biomimetic Synthetic Musculotendon Unit
%  Structural Material Discovery via Coherence Field Framework
%  Version 1.0 (DRAFT) — June 2026
%  Confidential — Internal Use Only
% ============================================================
\documentclass[11pt, a4paper]{article}

\usepackage[margin=2.2cm]{geometry}
\usepackage{amsmath, amssymb}
\usepackage{booktabs, array, longtable}
\usepackage{xcolor}
\usepackage{hyperref}
\usepackage{titlesec}
\usepackage{enumitem}
\usepackage{fancyhdr}
\usepackage{setspace}

\hypersetup{colorlinks=true, linkcolor=blue!50!black, urlcolor=blue!60!black}
\definecolor{nm-dark}{RGB}{30,30,30}
\definecolor{nm-gray}{RGB}{100,100,100}
\definecolor{nm-green}{RGB}{20,100,40}
\definecolor{nm-red}{RGB}{160,30,30}
\definecolor{nm-orange}{RGB}{180,100,20}

\pagestyle{fancy}\fancyhf{}
\rhead{\textcolor{nm-gray}{\small CFF Myofiber Analysis v1.0 (DRAFT)}}
\lhead{\textcolor{nm-gray}{\small Confidential --- Internal Use Only}}
\cfoot{\thepage}

\titleformat{\section}{\large\bfseries\color{nm-dark}}{\thesection}{0.6em}{}[\titlerule]
\titleformat{\subsection}{\normalsize\bfseries\color{nm-dark}}{\thesubsection}{0.6em}{}

\setlength{\parindent}{0pt}\setlength{\parskip}{6pt}\onehalfspacing

\begin{document}

\begin{center}
  {\LARGE \textbf{CFF Discovery Analysis}}\\[6pt]
  {\large Biomimetic Synthetic Musculotendon}\\[2pt]
  {\large for Substrate-Driven Embodiment}\\[10pt]
  {\color{nm-gray}\small Version 1.0 DRAFT \quad|\quad June 2026 \quad|\quad Confidential}\\[4pt]
  {\color{nm-gray}\small Joseph Forrester \quad|\quad DSF-AI Architecture}\\[4pt]
  {\color{nm-gray}\small (Eve draft, Joe to verify axiom mapping and thresholds)}
\end{center}

\vspace{1em}\hrule\vspace{1em}

% ============================================================
\section*{Executive Summary}

Applying the CFF sequential constraint elimination algorithm to
the target property \textbf{``biomimetic synthetic musculotendon
unit with biological-grade performance and home-deployment cycle
life,''} the forcing chain produces:

\begin{enumerate}[nosep]
  \item \textbf{The optimal mechanism is hydraulic-actuated
        braided contractile (McKibben-class), monolithic with
        tendon.} Forced by the conjunction of high force density,
        sub-50~ms response, and zero sliding interface required
        for $10^8$-cycle fatigue life. Not a choice --- the unique
        survivor after six filters.
  \item \textbf{The optimal bladder polymer is a self-healing
        thermoplastic polyurethane with Diels--Alder dynamic
        crosslinks.} Forced by the dispositional preservation
        axiom applied to cyclic strain. Standard elastomers (latex,
        medical silicone, SBR) eliminated on fatigue.
  \item \textbf{The optimal braid is UHMWPE at $20\degree$ initial
        angle, with elastomeric interphase eliminating
        braid--bladder abrasion.} Forced by geometric mode
        confinement and weak-coupling axioms.
\end{enumerate}

\textbf{This analysis takes the CFF algorithm developed for
superconductor discovery, applied unmodified in structure to a
soft-matter chemomechanical actuator domain.} The same six axioms
apply; only the databases change. The output is a manufacturable
recipe with documented forcing chain and rejected alternatives.

\tableofcontents

% ============================================================
\section{Target Property}

\begin{center}
\begin{tabular}{@{}ll@{}}
\toprule
\textbf{Target} & Biomimetic synthetic musculotendon unit \\
\textbf{Unloaded contraction} & $\geq 30\%$ of resting length \\
\textbf{Response time} & $\leq 50$~ms (full contraction) \\
\textbf{Force-to-weight} & $\geq 1$~kg force per 3~g fiber \\
\textbf{Cyclic operation} & 0.1--10~Hz biological range \\
\textbf{Cycle life} & $\geq 10^8$ cycles ($\sim$10~yr home service) \\
\textbf{Hysteresis} & $\leq 5\%$ at 30\% strain \\
\textbf{Working fluid} & Water-compatible (per Clone precedent) \\
\textbf{Operating pressure} & $\leq 1.0$~MPa (skin-contact safety) \\
\textbf{Manufacturability} & Continuous extrusion at meter-scale\\
\bottomrule
\end{tabular}
\end{center}

\textbf{Note on Clone benchmark:} Clone's published Myofiber
specifications ($\sim$1~kg / 3~g / $<$50~ms / $>$30\%) are taken
as the empirical existence proof that the target is physically
achievable. The CFF analysis derives the architectural class that
must satisfy these targets without assuming Clone's specific
chemistry.

% ============================================================
\section{Axiom-to-Soft-Matter Mapping}

The CFF axioms applied in the RTSC and Metamaterial discoveries
translate to this domain as follows:

\begin{center}
\begin{tabular}{@{}cp{4.5cm}p{8cm}@{}}
\toprule
\textbf{Axiom} & \textbf{Original sense} & \textbf{Soft-matter sense} \\
\midrule
A3 & Boundary conditions & Tendon attachment must not impose
    competing mechanical or fluidic modes \\
A5 & Dispositional preservation & Cyclic strain must not
    reorganize material structure (fatigue criterion) \\
A6 & Non-trivial winding / frustration & Geometric mechanism
    must localize / concentrate pressure-to-strain coupling \\
A7 & Coherence preservation (weak coupling) & Energy conversion
    must minimize dissipation: sliding friction, viscoelastic
    creep, hysteresis are loss channels \\
A8 & Local field response / coupling amplitude & Force per unit
    volume must exceed threshold for the kg-force target \\
\bottomrule
\end{tabular}
\end{center}

\textbf{Caveat:} This mapping is Eve's interpretation. Joe should
verify each translation. Specifically, A7 (``coherence
preservation'') applied to mechanical systems could equally be
read as ``minimize entropy production per cycle'' --- a stricter
criterion that would further constrain the bladder polymer class.
The looser reading is used below; the stricter reading should be
applied as a refinement once verified.

% ============================================================
\section{CFF Structural Filters}

\subsection{Filter 1: Topological Closure of Contractile Mechanism}

The actuator must convert distributed (3D) pressure or field
energy into directed (1D) length change without phase scrambling.
This requires a topological mechanism connecting input geometry
to output geometry.

\begin{center}
\begin{tabular}{@{}llrl@{}}
\toprule
\textbf{Mechanism} & \textbf{Closure} & \textbf{Strain ceiling} &
  \textbf{Survives?} \\
\midrule
Braided pneumatic / hydraulic & Braid angle & $\sim$30\% &
  \textcolor{nm-green}{\textbf{Yes}} \\
Twisted-coiled polymer (TCP) & Helical twist & $\sim$50\% &
  \textcolor{nm-orange}{Marginal (thermal speed)} \\
Liquid crystal elastomer (LCE) & Nematic order & $\sim$40\% &
  \textcolor{nm-orange}{Marginal (thermal speed)} \\
HASEL electrohydraulic & Incompressibility & $\sim$25\% &
  \textcolor{nm-orange}{Borderline strain} \\
Piezoelectric stack & Direct field & $<$0.5\% &
  \textcolor{nm-red}{No (strain)} \\
Shape memory alloy & Phase transition & $\sim$8\% &
  \textcolor{nm-red}{No (strain + speed)} \\
Dielectric elastomer (DEAP) & Coulomb pressure & $\sim$30\% &
  \textcolor{nm-red}{No (force density)} \\
Bulk hydrogel & Osmotic swelling & $>$50\% &
  \textcolor{nm-red}{No (diffusion-limited speed)} \\
\bottomrule
\end{tabular}
\end{center}

\textbf{Result:} \textcolor{nm-green}{Braided pneumatic/hydraulic
survives unambiguously.} TCP and LCE survive on strain but are
eliminated by Filter 4 (speed).

\subsection{Filter 2: Weak Coupling / Loss Minimization (A7)}

Coherent contraction requires energy to convert pressure
$\to$ strain with minimal dissipation. Loss channels in
contractile systems:

\begin{center}
\begin{tabular}{@{}lll@{}}
\toprule
\textbf{Loss channel} & \textbf{Consequence} & \textbf{Mitigation} \\
\midrule
Braid--bladder sliding friction & Fatigue, hysteresis, heat &
  Eliminate interface (monolithic) or lubricate \\
Viscoelastic polymer creep & Slow recovery, residual strain &
  Low-loss polymer chemistry \\
Bladder fold/wrinkle hysteresis & Energy loss per cycle &
  Pre-stress or geometric design \\
Fluid viscous dissipation & Speed limitation &
  Low-viscosity fluid + large ports \\
\bottomrule
\end{tabular}
\end{center}

\textbf{CFF filter:} Eliminate or strongly mitigate each loss
channel. The dominant channel for McKibben-class actuators at
$10^8$ cycles is braid--bladder abrasion. Two architectures
survive:

\begin{enumerate}[nosep]
  \item \textbf{Monolithic musculotendon}: bladder, braid analog,
        and tendon as a single polymer pour. No interface = no
        sliding wear. (Clone's published approach.)
  \item \textbf{Lubricated braid-bladder interface}: elastomeric
        interphase layer or PTFE liner between braid and bladder.
        Friction reduced by $\sim 10\times$, fatigue limit
        extended.
\end{enumerate}

\textbf{Result:} Both architectures survive Filter 2. Either
eliminates the primary fatigue mode.

\subsection{Filter 3: Geometric Mode Confinement (A6)}

The pressure-to-strain coupling must be \textit{localized} ---
distributed pressure must produce concentrated, directional
contraction. This is the analog of frustration / mode confinement
in the superconductor and metamaterial discoveries.

\begin{center}
\begin{tabular}{@{}lll@{}}
\toprule
\textbf{Geometry} & \textbf{Confinement mode} &
  \textbf{Strain coupling} \\
\midrule
Helical braid, $\alpha < 54.7\degree$ initial & Angle inversion &
  Strong (axial) \\
Helical braid, $\alpha > 54.7\degree$ initial & Angle inversion &
  Reverse (expansion) \\
Pleated bladder (Festo Fluidic Muscle) & Fold mechanics &
  Moderate \\
Series-elastic stretched braid & Tension transfer &
  Weak \\
Bulk elastomer pressure cell & None (isotropic) &
  None \\
\bottomrule
\end{tabular}
\end{center}

\textbf{CFF filter:} Helical braid at sub-magic-angle
($\alpha_0 \approx 20\degree$ optimum) maximizes axial coupling
per unit pressure. The magic angle $54.7\degree$ is the topological
inversion point where contraction switches to expansion --- a
physical constant from braid geometry, not a tuning parameter.

\textbf{Result:} Helical braid at $\alpha_0 = 20\degree \pm 2\degree$
forced. Other geometries are inferior or non-contractile.

\subsection{Filter 4: Force Density / Coupling Amplitude (A8)}

Force per unit volume must support the target: $\sim$1~kg force
from a 3~g fiber, roughly $\sim$3~cm$^3$ at polymer density.
That implies $\sim$110~kPa equivalent stress at the tendon.

\begin{center}
\begin{tabular}{@{}lrl@{}}
\toprule
\textbf{Mechanism @ working condition} & \textbf{Force density (MJ/m$^3$)} &
  \textbf{Survives?} \\
\midrule
Hydraulic McKibben @ 1.0~MPa & 3.0 &
  \textcolor{nm-green}{\textbf{Yes}} \\
Hydraulic McKibben @ 0.5~MPa & 1.5 &
  \textcolor{nm-green}{\textbf{Yes}} \\
Pneumatic McKibben @ 0.5~MPa & 1.5 &
  \textcolor{nm-green}{\textbf{Yes}} \\
TCP (peak, slow) & 3.0 &
  \textcolor{nm-orange}{Speed-limited} \\
LCE (thermal) & 0.5 &
  \textcolor{nm-red}{Below threshold} \\
HASEL & 0.3 &
  \textcolor{nm-red}{Below threshold} \\
Biological skeletal muscle & 0.3 &
  (reference) \\
\bottomrule
\end{tabular}
\end{center}

\textbf{Note:} Biological muscle operates at $\sim$0.3~MJ/m$^3$
because it is bandwidth-optimized, not force-optimized. The CFF
target (kg/3g = 10$\times$ biological) is unachievable with pure
biological mechanisms and is enabled in synthetic muscle by
higher working pressure.

\textbf{Result:} Hydraulic McKibben at 0.3--1.0~MPa survives with
margin. Above 1.0~MPa, fatigue limits dominate (Filter 5).

\subsection{Filter 5: Speed / Bandwidth}

Response time $\leq 50$~ms requires:

\begin{itemize}[nosep]
  \item Bladder wall stiffness sufficient to respond elastically
        (not viscoelastically) at the required strain rate
  \item Fluid inlet area large relative to bladder volume
        (low fluidic time constant)
  \item Working fluid with low kinematic viscosity
  \item No phase transition or diffusion-limited step in the
        mechanism
\end{itemize}

\begin{center}
\begin{tabular}{@{}llrl@{}}
\toprule
\textbf{Mechanism} & \textbf{Limit} & \textbf{Typical response} &
  \textbf{Survives?} \\
\midrule
Hydraulic McKibben (water) & Fluidic & 10--40~ms &
  \textcolor{nm-green}{\textbf{Yes}} \\
Pneumatic McKibben (air) & Compressibility & 30--80~ms &
  \textcolor{nm-orange}{Marginal} \\
TCP (thermal) & Heat diffusion & 1--10~s &
  \textcolor{nm-red}{No} \\
LCE (thermal) & Heat diffusion & 0.5--5~s &
  \textcolor{nm-red}{No} \\
SMA (thermal) & Heat diffusion & 0.5--3~s &
  \textcolor{nm-red}{No} \\
\bottomrule
\end{tabular}
\end{center}

\textbf{Result:} Hydraulic survives cleanly. Pneumatic is
marginal --- adequate for development and proof-of-concept, but
hydraulic forced for the production target.

\textbf{This is a non-obvious conclusion.} Pneumatic actuation is
cheaper and easier to plumb (compressed air, no fluid lines, no
leaks-stain-the-floor failure mode). The CFF chain forces
hydraulic because pneumatic's compressibility introduces
bandwidth loss that violates the $\leq 50$~ms requirement at
production-scale muscle counts. \textbf{This explains Clone's
water choice as a forcing-chain consequence, not a design
preference.}

\subsection{Filter 6: Dispositional Preservation under Cycling (A5)}

Material structure must not reorganize under cyclic loading. The
fatigue threshold for home deployment ($\sim$10 yr at 1 Hz duty
cycle) is approximately $10^8$ cycles.

\begin{center}
\begin{tabular}{@{}lrl@{}}
\toprule
\textbf{Bladder polymer class} & \textbf{Fatigue (cycles to failure
  at 30\% strain)} & \textbf{Survives?} \\
\midrule
Natural latex rubber & $10^4$--$10^5$ &
  \textcolor{nm-red}{No} \\
Medical silicone (LSR) & $10^5$--$10^6$ &
  \textcolor{nm-red}{No} \\
SBR / NBR & $10^5$--$10^6$ &
  \textcolor{nm-red}{No (high hysteresis)} \\
Standard TPU (Pellethane class) & $10^6$--$10^7$ &
  \textcolor{nm-orange}{Marginal} \\
LCE network & $10^7$ &
  \textcolor{nm-orange}{Marginal} \\
Self-healing TPU (Diels--Alder) & $10^8$--$10^9$ &
  \textcolor{nm-green}{\textbf{Yes}} \\
Self-healing TPU (disulfide exchange) & $10^8$+ &
  \textcolor{nm-green}{\textbf{Yes}} \\
Vitrimer (transesterification) elastomer & $10^8$+ &
  \textcolor{nm-green}{\textbf{Yes}} \\
\bottomrule
\end{tabular}
\end{center}

\textbf{CFF filter:} Standard polymers fail at the $10^8$ threshold
by 1--3 orders of magnitude. Dynamic covalent network polymers
survive because they continuously repair fatigue microcracks ---
healing rate exceeds damage accumulation rate at biological cycle
rates.

\textbf{Result:} Self-healing dynamic-covalent network elastomers
forced. Three chemistries survive (see Section 5).

\subsection{Filter 7: Tendon Boundary Condition (A3)}

The muscle-to-bone attachment is where most McKibben failures
historically occur. Stress concentrates at the termination,
fatigue cracks initiate at the bladder-clamp interface.

\begin{center}
\begin{tabular}{@{}ll@{}}
\toprule
\textbf{Termination architecture} & \textbf{Survives?} \\
\midrule
Bolted clamp on cylindrical end & \textcolor{nm-red}{No (stress concentration)} \\
Knotted braid termination & \textcolor{nm-red}{No (abrasion)} \\
Crimp + strain relief & \textcolor{nm-orange}{Marginal} \\
Adhesive bond, modulus-matched & \textcolor{nm-orange}{Marginal} \\
Monolithic musculotendon (continuous extrusion) &
  \textcolor{nm-green}{\textbf{Yes}} \\
\bottomrule
\end{tabular}
\end{center}

\textbf{CFF filter:} The termination must not introduce a
modulus discontinuity or stress concentration. Monolithic
construction --- where bladder, braid analog, and tendon are
formed in one continuous polymer process with graded crosslink
density --- eliminates the interface entirely.

\textbf{Result:} Monolithic musculotendon forced. This is the
specific feature Clone has emphasized publicly, and the CFF
chain independently arrives at the same constraint.

% ============================================================
\section{Forcing Chain}

\begin{enumerate}
  \item \textbf{Strain target $\geq 30\%$} (Filter 1) $\to$
        eliminates piezo, SMA, dielectric elastomer

  \item \textbf{Response $\leq 50$~ms} (Filter 5) $\to$
        eliminates TCP, LCE, SMA, marginalizes pneumatic
        $\to$ \textbf{hydraulic forced}

  \item \textbf{Force density $\geq 1$~MJ/m$^3$} (Filter 4) $\to$
        confirms McKibben class, eliminates LCE / HASEL

  \item \textbf{Mode confinement} (Filter 3) $\to$
        helical braid at $\alpha_0 \approx 20\degree$ forced
        \textit{(magic-angle topology, not tuned)}

  \item \textbf{Loss minimization} (Filter 2) $\to$
        monolithic construction OR lubricated interface
        $\to$ \textbf{monolithic preferred for cycle life}

  \item \textbf{Fatigue $\geq 10^8$ cycles} (Filter 6) $\to$
        \textbf{dynamic-covalent self-healing elastomer forced}
        $\to$ eliminates latex, silicone, standard TPU

  \item \textbf{Termination boundary} (Filter 7) $\to$
        \textbf{monolithic musculotendon forced} (re-confirms F2)

  \item \textbf{Working fluid} (residual constraint) $\to$
        water-based, low viscosity, biocompatible, thermal
        stability $\to$ \textbf{water + propylene glycol +
        corrosion inhibitor}
\end{enumerate}

The chain is deterministic. Given the same materials databases and
target thresholds, any implementation produces the same
architectural class.

% ============================================================
\section{Architectural Class (Output)}

\begin{center}
\begin{tabular}{@{}lp{9cm}@{}}
\toprule
\textbf{Element} & \textbf{Forced value} \\
\midrule
Actuation mechanism & Hydraulic-pressurized helical braided bladder
  (McKibben-class) \\
Operating pressure & 0.3--0.8~MPa nominal, 1.2~MPa burst \\
Working fluid & Distilled H$_2$O + 5--10\% propylene glycol +
  corrosion inhibitor \\
Bladder polymer & Self-healing dynamic-covalent elastomer
  (Diels--Alder, disulfide-exchange, or vitrimer class) \\
Bladder modulus & 50--200~kPa (Shore A 60--80 equivalent) \\
Bladder hysteresis & $\leq 5\%$ at 30\% strain \\
Braid material & UHMWPE filament (alt: Vectran for higher
  modulus) \\
Braid initial angle & $20\degree \pm 2\degree$ from axis \\
Braid--bladder interface & Elastomeric interphase OR fully
  monolithic (preferred) \\
Tendon material & Same elastomer chemistry, graded crosslink
  density to $\sim$10~MPa effective modulus over 2--3~mm
  transition \\
Termination & Monolithic musculotendon (continuous extrusion),
  no mechanical interface \\
Geometry & Length-to-diameter ratio $\sim$10:1 \\
\bottomrule
\end{tabular}
\end{center}

% ============================================================
\section{Ranked Candidate Chemistries}

The bladder polymer is the highest-information variable in the
architectural class. Three chemistries survive all filters; they
are ranked by manufacturability and published validation.

\subsection{Rank 1: Furan--Maleimide Diels--Alder TPU}

\begin{center}
\begin{tabular}{@{}lp{9cm}@{}}
\toprule
\textbf{Property} & \textbf{Value} \\
\midrule
Network chemistry & TPU backbone with furan and maleimide
  pendant groups forming reversible Diels--Alder crosslinks \\
Healing temperature & $\sim$80--90$\degree$C (cycle:
  retro-DA $\to$ cool $\to$ re-DA) \\
Operating range & Room temp to 60$\degree$C (network stable) \\
Demonstrated fatigue & $10^9$+ cycles in published cyclic tests
  at 50\% strain \\
Hysteresis & $<$3\% at 30\% strain \\
Modulus (range) & 50--500~kPa tunable via crosslink density \\
Synthesis & Published (Chen et al., 2002 and successors).
  Standard TPU polycondensation with bismaleimide chain extender
  and furfuryl-functional diol. \\
\midrule
Manufacturing risk & \textcolor{nm-green}{Low.} Components
  commercially available; chemistry runs in standard PU reactor;
  monolithic extrusion compatible. \\
Healing trigger concern & \textcolor{nm-orange}{Body-temp
  operation is below DA healing temperature. Healing requires
  periodic thermal cycling (overnight rest at $\sim$50$\degree$C)
  or in-service warming via fluid temperature.} \\
\bottomrule
\end{tabular}
\end{center}

\subsection{Rank 2: Disulfide-Exchange Polyurethane}

\begin{center}
\begin{tabular}{@{}lp{9cm}@{}}
\toprule
\textbf{Property} & \textbf{Value} \\
\midrule
Network chemistry & TPU with aromatic disulfide crosslinks that
  exchange continuously at room temperature \\
Healing temperature & Room temp (continuous, no trigger needed) \\
Operating range & --20 to 80$\degree$C \\
Demonstrated fatigue & $10^8$+ cycles at 30\% strain \\
Hysteresis & $\sim$5--8\% at 30\% strain
  \textcolor{nm-orange}{(higher than Rank 1)} \\
Modulus & 100--300~kPa typical \\
Synthesis & Published (Rekondo et al., 2014). Commercial
  precursors (aromatic disulfide diols) available.\\
\midrule
Manufacturing risk & \textcolor{nm-green}{Low--medium.} Slightly
  more complex monomer synthesis than Rank 1. \\
Operational advantage & \textcolor{nm-green}{Continuous
  self-healing at body / room temperature --- no thermal cycling
  required for service life.} \\
\bottomrule
\end{tabular}
\end{center}

\subsection{Rank 3: Transesterification Vitrimer Elastomer}

\begin{center}
\begin{tabular}{@{}lp{9cm}@{}}
\toprule
\textbf{Property} & \textbf{Value} \\
\midrule
Network chemistry & Topology-preserving covalent network with
  catalyzed transesterification exchange \\
Healing temperature & 120--160$\degree$C (well above operating) \\
Operating range & Room temp to 100$\degree$C \\
Demonstrated fatigue & $10^8$ cycles, lower hysteresis than DA \\
Hysteresis & $\sim$2--4\% at 30\% strain \\
Modulus & Highly tunable (10~kPa to 10~MPa) \\
\midrule
Manufacturing risk & \textcolor{nm-orange}{Medium.} Catalyst
  selection and dispersion non-trivial; higher cure temperature
  complicates monolithic extrusion of bladder + tendon graded
  network. \\
Strategic value & \textcolor{nm-green}{Best fatigue / hysteresis
  combination known; reserve for v2 production line once basic
  process is proven.} \\
\bottomrule
\end{tabular}
\end{center}

% ============================================================
\section{Rejected Families}

\begin{center}
\begin{tabular}{@{}lp{8cm}@{}}
\toprule
\textbf{Family} & \textbf{Reason for rejection} \\
\midrule
Natural latex rubber & Fatigue $<10^6$ at 30\% strain \\
Medical silicone (Pt-cure LSR) & Fatigue $<10^7$; surface
  abrasion under braid contact \\
SBR / NBR / EPDM & High hysteresis $\to$ thermal buildup and
  cycle loss \\
Standard thermoset polyurethane & No self-healing pathway $\to$
  fails $10^8$ threshold \\
Liquid crystal elastomers (thermal) & Diffusion-limited speed
  $\to$ fails 50~ms threshold \\
Twisted-coiled polymers (TCP) & Heat-driven contraction $\to$
  fails 50~ms threshold; otherwise excellent force density \\
Shape memory alloys & Speed limit + strain limit \\
HASEL electrohydraulic & Force density below threshold \\
Dielectric elastomers (DEAP) & Force density below threshold;
  high-voltage driver inappropriate for embodied use \\
Pneumatic McKibben & Compressibility violates 50~ms response
  at production scale \textit{(retained for development /
  prototyping only)} \\
\bottomrule
\end{tabular}
\end{center}

% ============================================================
\section{Manufacturing Specification}

\subsection{Continuous Monolithic Musculotendon Process}

\begin{enumerate}[nosep]
  \item Synthesize self-healing TPU pre-polymer (Rank 1 or 2
        chemistry). Maintain inventory of two stock streams:
        ``soft'' (low crosslink density, bladder-grade,
        $\sim$100~kPa modulus) and ``stiff'' (high crosslink
        density, tendon-grade, $\sim$5--10~MPa modulus).
  \item Co-extrude through a profile die that grades from soft
        to stiff along axial position, producing a continuous
        rod with bladder section, tapered transition zone, and
        tendon ends.
  \item In a second station, integrate UHMWPE filament via
        rotary braiding around the bladder section as it
        emerges, at the specified $20\degree$ angle.
        \textit{Alternative monolithic path: embed continuous
        UHMWPE strands during co-extrusion, eliminating the
        braid station entirely. This is the path Clone is most
        likely using based on the ``produced by the kilometer''
        manufacturing claim.}
  \item Cure / set the network in line (DA: thermal cure;
        disulfide: ambient set).
  \item Cap ends with valve fittings (or in-line crimp the
        tendons to mounting hardware). Cap interfaces use the
        same polymer chemistry to maintain continuity.
  \item Pressure-test each unit to 1.5$\times$ working pressure;
        record burst sample data per lot.
\end{enumerate}

\subsection{Minimum Viable v0 Recipe (for First Test Batches)}

Before scaling to monolithic continuous extrusion, a hand-buildable
proof-of-architecture can be assembled from:

\begin{itemize}[nosep]
  \item \textbf{Bladder tube}: Rank 2 (disulfide-exchange TPU),
        compression-molded as a 8--12~mm OD tube, 1~mm wall,
        room-temperature self-healing
  \item \textbf{Braid}: 16-carrier UHMWPE braid, $20\degree$
        initial angle, lubricated with silicone oil at braid-bladder
        interface (Filter 2 alternative path)
  \item \textbf{End caps}: 3D-printed PETG with O-ring seal,
        bonded with TPU-compatible adhesive (DCM-activated)
  \item \textbf{Working fluid}: distilled water + 10\% propylene
        glycol + 0.1\% benzotriazole corrosion inhibitor
  \item \textbf{Operating pressure}: start at 0.3~MPa for
        characterization, ramp to 0.6~MPa for performance test
\end{itemize}

This v0 recipe lets us validate force-strain curve and bandwidth
on the bench while the production process is engineered.

% ============================================================
\section{Manufacturing Partners}

For prototype-quantity production of the bladder polymer:

\begin{itemize}[nosep]
  \item \textbf{Specialty TPU custom synthesis}: Lubrizol
        (Pellethane / Estane variants), Covestro, BASF Elastollan
        custom group. Each has done custom self-healing TPU
        formulations for medical/research customers.
  \item \textbf{Polymer chemistry contract synthesis}:
        Polymer Source (Montreal), Sigma-Aldrich custom synthesis,
        Strem custom polymers. For Rank 1 (DA-TPU) chemistry,
        Polymer Source has published precedent.
  \item \textbf{Co-extrusion services}: Plastic Extrusion
        Technologies, Polymer Extrusion Technology Inc. for
        prototype runs; Microspec Corp for medical-grade
        precision profiles.
  \item \textbf{Braiding}: Steeger USA, Wardwell Braiding
        for industrial-precision braiders capable of the
        $20\degree \pm 2\degree$ tolerance.
\end{itemize}

Recommended first contact: a contract polymer chemistry lab to
produce 1--2~kg of Rank 2 (disulfide-exchange TPU) stock for
manual bladder fabrication, in parallel with a co-extrusion
quote for the monolithic v1 path.

% ============================================================
\section{What Makes This Defensible}

\subsection{Auditable Inputs}

\begin{itemize}[nosep]
  \item Polymer fatigue data from published cyclic-loading studies
  \item McKibben kinematics from established braid mechanics
        (Chou \& Hannaford, Tondu, etc.)
  \item Polymer chemistry references with citations
\end{itemize}

\subsection{Reproducible Output}

The architectural class is deterministic given the filters and
thresholds. A reviewer who disagrees with a threshold can re-run
with a different value and see which candidates are gained or lost.

\subsection{Honest Scope}

The CFF chain forces an architectural class. Specific composition
parameters (crosslink density, exact braid pitch, fluid additive
ratios) are tuning parameters that require empirical validation.
The chain says \textit{what kind of muscle}; benchwork says
\textit{which exact recipe within that kind}.

\subsection{Independent Convergence with Clone}

The chain independently arrives at:

\begin{itemize}[nosep]
  \item Water-driven actuation (Filter 5: hydraulic forced over
        pneumatic)
  \item Monolithic musculotendon (Filters 2 + 7: termination
        + loss minimization)
  \item Polymer rather than rubber bladder (Filter 6: fatigue)
\end{itemize}

These match Clone's publicly stated design choices. CFF predicts
them as forced consequences, not design preferences. This is
either (a)~CFF correctly identifying the unique solution Clone
also found, or (b)~Clone's choices and CFF's outputs both
reflecting the same underlying physics. Either reading supports
the architectural class.

\subsection{Departure from Clone}

CFF does \textit{not} force Clone's specific polymer chemistry.
Their proprietary formulation is one member of the architectural
class. Other members (Rank 2, Rank 3) are equally valid solutions
to the forced class. \textbf{The architectural class is the moat,
not the specific polymer.} A manufacturable Myofiber-equivalent
can be produced without infringing Clone's patents by using a
different point in the same class.

% ============================================================
\section{Open Questions for Verification}

\begin{enumerate}[nosep]
  \item \textbf{Axiom translation.} Joe to verify that the
        soft-matter mapping of A3, A5, A6, A7, A8 (Section 2)
        matches the formal CFF definitions. Especially A7's
        weak-coupling sense in mechanical systems.
  \item \textbf{Healing threshold.} The $10^8$ cycle target may
        be too aggressive (or insufficiently aggressive) depending
        on actual home duty cycle. Should be derived from
        embodiment use case.
  \item \textbf{Stricter A7 reading.} A stricter ``minimize entropy
        production per cycle'' interpretation may further narrow
        the candidate polymer chemistries. Worth re-running.
  \item \textbf{Hydration / swelling.} Water-contact for the
        bladder polymer raises hydrolysis questions. PU is
        generally hydrolysis-resistant in disulfide-exchange
        formulations; explicit confirmation needed for service life.
  \item \textbf{Sensor integration.} The architectural class does
        not yet incorporate the proprioceptive sensor requirement
        (muscle length + pressure). A v2 of this analysis should
        add a sensor-coupling filter.
\end{enumerate}

% ============================================================
\section{Recommended Next Steps}

\begin{enumerate}[nosep]
  \item \textbf{Verify axiom mapping (Section 2)} --- Joe, 30~min
        review.
  \item \textbf{Build v0 hand muscle} per Section 8.2 spec.
        Single unit, instrument force/strain/pressure, validate
        50~ms response and 30\% strain. Bench validation of the
        architecture before committing to production process.
  \item \textbf{Contract polymer chemistry lab} for 1--2~kg
        Rank 2 (disulfide-exchange TPU) stock. Use for v0 muscle.
  \item \textbf{Parallel}: co-extrusion quote for monolithic v1
        path. Lead time to first production sample $\sim$3--6~mo.
  \item \textbf{v1 validation}: 4--8 monolithic muscles, full
        characterization including $10^7$+ cycle fatigue test.
        \textit{Note: $10^8$ test requires $\sim$3.2 years at 1~Hz;
        accelerated test at higher frequency standard.}
  \item \textbf{Strategic IP review}: the architectural class is
        defensible; specific recipes within the class may be
        patentable distinct from Clone's chemistry.
\end{enumerate}

% ============================================================
\section{Version History}

\begin{center}
\begin{tabular}{@{}lp{10.5cm}@{}}
\toprule
\textbf{Version} & \textbf{Description} \\
\midrule
1.0 DRAFT (June 2026) & Initial application of CFF Discovery
  Algorithm to biomimetic synthetic musculotendon. Seven structural
  filters, deterministic forcing chain, architectural class
  defined, three polymer chemistries ranked, manufacturing path
  to v0 and v1 specified. Independent convergence with Clone's
  public design choices noted. Open questions flagged for Joe
  verification, especially axiom-to-soft-matter mapping. \\
\bottomrule
\end{tabular}
\end{center}

\vspace{2em}
\hrule
\vspace{0.5em}
{\color{nm-gray}\small This document is confidential and for
internal use only. The CFF Discovery Algorithm is the intellectual
property of Joseph Forrester / DSF-AI. The architectural class
output for synthetic musculotendon is novel application of the
existing framework. The underlying UFCP framework is a trade
secret. The polymer chemistry references use published values
and are independently verifiable.}

\end{document}
